\documentclass[aps,prd,reprint,superscriptaddress,nofootinbib,amsmath,amssymb,floatfix,longbibliography]{revtex4-2}
\usepackage[T1]{fontenc}
\usepackage{booktabs}
\usepackage{graphicx}
\usepackage{xurl}
\usepackage[colorlinks=true,linkcolor=blue,citecolor=blue,urlcolor=blue]{hyperref}

\newcommand{\manuscripttitle}{Emulating Cosmic Reionization Histories and 21-cm Power Spectra for Nonstandard Small-Scale Matter Power}
\hypersetup{pdftitle={\manuscripttitle},pdfauthor={Haina Huang}}

\setlength{\emergencystretch}{2em}
\renewcommand{\textfraction}{0.08}
\renewcommand{\topfraction}{0.92}
\renewcommand{\bottomfraction}{0.92}
\renewcommand{\floatpagefraction}{0.80}
\renewcommand{\dbltopfraction}{0.92}
\renewcommand{\dblfloatpagefraction}{0.80}
\setlength{\dblfloatsep}{8pt plus 2pt minus 2pt}
\setlength{\dbltextfloatsep}{10pt plus 2pt minus 2pt}
\setcounter{topnumber}{4}
\setcounter{dbltopnumber}{1}
\setcounter{totalnumber}{6}

\newcommand{\figwide}[3]{%
  \begin{figure*}[!tb]
    \centering
    \includegraphics[width=168mm]{#1}
    \caption{#2}\label{#3}
  \end{figure*}%
}

\newcommand{\figcolumn}[3]{%
  \begin{figure}[!tbp]
    \centering
    \includegraphics[width=\columnwidth]{#1}
    \caption{#2}\label{#3}
  \end{figure}%
}

\begin{document}
\title{\manuscripttitle}
\author{Haina Huang}
\affiliation{Department of Physics, The Chinese University of Hong Kong, Shatin, Hong Kong, China}
\date{November 26, 2026}

\begin{abstract}
Jointly constraining astrophysical galaxy formation parameters and enhancements in the small-scale linear matter power spectrum requires repeated cosmological reionization and 21-cm power-spectrum evaluations across eight to ten parameters, rendering direct seminumerical Bayesian sampling computationally prohibitive. We present a unified ten-dimensional emulation framework for a broken-power-law matter spectrum integrated into \textsc{21cmFAST}, comprising a standalone global neutral-fraction history emulator, full-box simulation data engineering (\texttt{ps\_xhi\_fullbox\_v1}), the joint 21-cm power-spectrum and history surrogate (\textsc{PSX-emu v1.0}), and localized residual patching over ten parameters ($k_{p,h} \in [1, 30]\,h\,\mathrm{Mpc}^{-1}$, $m_s \in [0.5, 4.0]$, and eight astrophysical coordinates). The history emulator averages five pre-normalized residual multilayer perceptrons ($1{,}065{,}248$ parameters per member; $\sim 0.41\,\mathrm{ms}$ five-seed evaluation versus $416\,\mathrm{s}$ median direct box wall time) trained on 100,000 full-box runs ($250\,\mathrm{Mpc}$, $512^3/128^3$, pinned at \texttt{N\_THREADS=16} because GNU Scientific Library Fourier-mode initialization depends on OpenMP thread count). Scaling from 65,536 to 100,000 runs reduces 512-history development RMSE in $x_{\mathrm{HI}}(5.9)$ by $35.6\%$ ($0.00699 \to 0.00450$), $\tau$ RMSE by $23.2\%$ ($8.74\times 10^{-4} \to 6.72\times 10^{-4}$), and $q_{90}(|\Delta\ell_{\mathrm{hist}}|)$ by $26.6\%$ ($1.612 \to 1.183$), while an out-of-sample 56-run posterior validation confirms $\ell_{\mathrm{reion}}$ $q_{90}$ errors of $0.056$ (8D) and $0.026$ (6D) and eight four-IC checks match exact four-IC marginal likelihoods within $\le 0.028$. Controlled power-spectrum pilots ($m_s \in [1.51, 4.0]$) show that 1-IC parameter coverage dominates 4-IC replication ($0.581$ vs.\ $0.579\,\mathrm{dex}$ at $N=534$; $0.344\,\mathrm{dex}$ at $N=2{,}326$) and expose parameter-dependent $1.48\times$ low-$k$ bias in 12-subbox averaging. On the completed 100,000-box \texttt{ps\_xhi\_fullbox\_v1} archive, the eight-seed multi-task \textsc{PSX-emu v1.0} surrogate ($878{,}209$ parameters per member, $\lambda=1.0$) resolves a $54/64$ train-validation split overlap in prior checkpoints and achieves $0.0485\,\mathrm{dex}$ overall $\log_{10}(\Delta^2_{21}+0.01)$ RMSE ($0.0334\,\mathrm{dex}$ on LF-posterior draws; $2.47\%$ median relative error at $\Delta^2_{21}>10\,\mathrm{mK}^2$) and $\tau$ RMSE $1.72\times 10^{-4}$ ($2.45\%$ of $\sigma_{\rm Planck}$) across 300 strictly disjoint independent points, under $2\times$ the $0.0250\,\mathrm{dex}$ two-IC stochastic floor. Finally, without retraining base weights, a 64-family GPU-accelerated localized radial-basis-function (RBF) residual patch ($512/128/250\,\mathrm{Mpc}$, $\texttt{ZPRIME\_STEP\_FACTOR}=1.005$) cuts held-out normalized prediction error $L_2\ p_{90}$ from $0.7924$ to $0.0671$ and $t_\star$ Fisher $L_3$ median from $0.973$ to $0.052$, removing the frozen-base $t_\star$ PS+LF posterior bias ($0.6326 \to 0.4942$, $0.524\sigma \to 0.017\sigma$ around truth $0.500$). Extending this local surrogate to anisotropic Mat\'ern-5/2 automatic-relevance-determination Gaussian-process (ARD-GP) patches across four lightcone/coeval benchmarks suppresses held-out $L_2\ p_{90}$ residuals to $0.0278\text{--}0.0656$ (from $0.6699\text{--}3.9249$), enabling three-way 8D hybrid Stretch+GMM MCMC chains ($230{,}400$ post-burn samples each) at $\mathrm{PL}\,(10,0.968)$ ($\max\widehat{R}=1.0071$, $|\mathrm{med\_z}|\le 0.383$), $\mathrm{BT}\,(10,2.5)$ ($\max\widehat{R}=1.0105$), and $\mathrm{BT}\,(3,2.0)$ ($\max\widehat{R}=1.0049$).
\end{abstract}

\maketitle

\section{Introduction}\label{sec:intro}

High-redshift ultraviolet (UV) galaxy luminosity functions (LFs), cosmic microwave background (CMB) optical depth measurements, Ly$\alpha$ forest dark-pixel bounds, and interferometric 21-cm power-spectrum upper limits probe overlapping astrophysical and cosmological physics during cosmic dawn and the epoch of reionization (EoR) \cite{Mesinger2011,Park2019,Planck2018,McGreer2015,Donnan2024,HERA2023}. When nonstandard small-scale primordial or dark-matter physics enhances the matter power spectrum above wavenumbers $k \gtrsim 1\,h\,\mathrm{Mpc}^{-1}$ \cite{Parashari23JWSTPPS,Hira24bluetilt,MunozDvorkin2020}, the resulting excess of low-mass halos shifts the halo mass--UV magnitude relation and reionization timing. Constraining small-scale spectral parameters robustly therefore requires marginalizing simultaneously over flexible galaxy-formation, ionizing-escape, and intergalactic medium (IGM) heating prescriptions rather than holding astrophysics fixed. While UV LFs can be evaluated rapidly from extended Press--Schechter halo mass functions, three-dimensional seminumerical reionization and 21-cm simulations in \textsc{21cmFAST} \cite{Mesinger2011,Murray2020} remain far too costly per evaluation for exhaustive multi-prior Bayesian inference across eight to ten dimensions without surrogate acceleration.

Neural emulation of \textsc{21cmFAST} summary observables is a well-established methodology. Notably, \textsc{21cmEMU} \cite{Breitman2024} demonstrated fast accurate surrogates for standard $\Lambda$CDM astrophysics, and \textsc{21cmEMUv3} \cite{Breitman2026} extended emulation with hybrid diffusion--long short-term memory (LSTM) architectures incorporating explicit initial-condition (IC) stochasticity. The ten-dimensional reionization emulation framework described in this work addresses a distinct, complementary regime: joint emulation across an extended ten-dimensional domain combining the eight astrophysical parameters of Park \textit{et al.}~\cite{Park2019} with two phenomenological parameters characterizing a broken small-scale linear matter power spectrum---break scale $k_{p,h}$ and small-scale tilt $m_s$---evaluated under frozen production box, resolution, physics, and OpenMP threading contracts.

The present methods paper deliberately isolates surrogate architecture, frozen simulation data contracts, numerical reproducibility, random-IC likelihood target semantics, and data-engineering lessons for full-box 21-cm power-spectrum production from the astrophysical constraint analysis reported in our companion physical inference paper, which deploys the validated 10D history emulator strictly as a calibrated reionization likelihood component alongside exact UV LFs. Where exploratory 21-cm power-spectrum pilots span a narrower slope range or simplified estimator, we explicitly disclose both domain boundaries and estimator differences so that every reported metric carries unambiguous provenance and out-of-sample validity.

The remainder of this paper is organized as follows. Section~\ref{sec:model} specifies the modified matter power spectrum, ten-dimensional design support, simulation data contracts (\texttt{ps\_xhi\_fullbox\_v1}), the discovery of GNU Scientific Library (GSL) OpenMP thread dependence in \textsc{21cmFAST} IC generation, and exact likelihood interfaces. Section~\ref{sec:history} details the history-emulator architecture, five-seed ensembling, paired 65,536-to-100,000 accuracy scaling on development-validation splits, and a finite smoke-chain comparison against exact simulations. Section~\ref{sec:ic} formalizes representative-history versus IC-marginalized likelihood targets, reconciling prior-box IC sensitivity against out-of-sample stratified direct-simulation validation inside posterior support. Section~\ref{sec:ps} presents controlled 21-cm power-spectrum data-engineering pilots (training-budget scaling and subbox-versus-full-box estimator bias) alongside the completed 100,000-box full-box production contract. Section~\ref{sec:psx_emu} introduces the production multi-task \textsc{PSX-emu v1.0} joint surrogate, documents our 300-point zero-overlap independent audit and four-way PL convergence diagnostics, validates localized RBF and anisotropic Mat\'ern-5/2 ARD-GP residual patches across coeval and lightcone mocks, and benchmarks three-way 8D hybrid Stretch+GMM posterior chains. Section~\ref{sec:discussion} discusses domain boundaries, future multi-ensemble MCMC validation, and runtime accounting, and Section~\ref{sec:conclusions} summarizes conclusions and code/data availability.

\section{Forward Simulation Model and Data Contracts}\label{sec:model}

\subsection{Broken-Power-Law Matter Spectrum and Galaxy Model}\label{sec:bt_model}

We modify the linear matter power spectrum within \textsc{21cmFAST} (\texttt{POWER\_SPECTRUM=6}) via a continuous broken power law normalized individually to fixed $\sigma_8 = 0.815$ alongside flat Planck 2015 cosmological parameters ($\Omega_m = 0.308$, $\Omega_b = 0.0484$, $h = 0.678$, $n_s = 0.968$) \cite{Planck2015,Eisenstein1998}:
\begin{equation}
P_{\mathrm{lin}}(k) = A(\sigma_8)\,T_{\mathrm{EH}}^2(k)
\begin{cases}
k^{n_s}, & k \le k_p,\\
k^{m_s}\,k_p^{\,n_s - m_s}, & k > k_p,
\end{cases}
\label{eq:pk_bt}
\end{equation}
where $T_{\mathrm{EH}}(k)$ is the Eisenstein--Hu transfer function \cite{Eisenstein1998}, $k_{p,h} \in [1, 30]\,h\,\mathrm{Mpc}^{-1}$ is converted internally to inverse-Mpc units via $k_{p,\mathrm{Mpc}^{-1}} = h\,k_{p,h}$ (i.e., $k_p$ in Eq.~\eqref{eq:pk_bt} is expressed in the same inverse-length convention as $k$), and $m_s \in [0.5, 4.0]$ (stored as \texttt{MS}). Holding $\sigma_8$ fixed adjusts the overall amplitude $A$ whenever enhanced small-scale power for $m_s > n_s$ contributes non-negligibly to the $8\,h^{-1}\,\mathrm{Mpc}$ variance. At $m_s = n_s = 0.968$, Eq.~\eqref{eq:pk_bt} reduces identically to the standard power-law (PL, \texttt{POWER\_SPECTRUM=0}) transfer branch \cite{Murray2020}.

Galaxy star formation and ionizing escape follow the halo-mass-dependent parameterization of Ref.~\cite{Park2019}:
\begin{align}
f_\star(M_h) &= \min\!\left[1,\, f_{\star,10}\left(\frac{M_h}{10^{10}\,M_\odot}\right)^{\alpha_\star}\right],\\
\dot{M}_\star &= \frac{f_\star(M_h)\,(\Omega_b/\Omega_m)\,M_h\,H(z)}{t_\star},\\
f_{\mathrm{esc}}(M_h) &= \min\!\left[1,\, f_{\mathrm{esc},10}\left(\frac{M_h}{10^{10}\,M_\odot}\right)^{\alpha_{\mathrm{esc}}}\right],\\
f_{\mathrm{duty}}(M_h) &= \exp\!\left(-\frac{M_{\mathrm{turn}}}{M_h}\right),
\end{align}
augmented by the soft X-ray luminosity per unit star formation rate $\log_{10} L_X/\mathrm{SFR}$ (\texttt{L\_X}) and photon energy threshold $E_0$ (\texttt{NU\_X\_THRESH}). Note that UV LFs depend on $\eta_\star \equiv \log_{10} f_{\star,10} - \log_{10} t_\star$, whereas the emulator input vector retains the orthogonal canonical parameterization $( \log_{10} f_{\star,10}, t_\star )$. Table~\ref{tab:params} summarizes the full ten-dimensional design support and frozen simulation flags.

\begin{table}[!tbp]
\caption{Ten-dimensional reionization-emulator parameter design support and frozen simulation contract parameters (\texttt{ps\_xhi\_fullbox\_v1}). All astrophysical parameters follow Ref.~\cite{Park2019}; the native binary is verified against SHA-256 digest \texttt{3211a6629109da73...94e5702b2}.}
\label{tab:params}
\footnotesize
\setlength{\tabcolsep}{3.5pt}
\begin{ruledtabular}
\begin{tabular}{llp{4.05cm}}
Coordinate & Range & Physical meaning / unit \\
\hline
\texttt{F\_STAR10} & $[-3.0,\, 0.0]$ & $\log_{10} f_{\star,10}$ (SF norm.) \\
\texttt{ALPHA\_STAR} & $[-0.5,\, 1.0]$ & $\alpha_\star$ (stellar mass slope) \\
\texttt{F\_ESC10} & $[-3.0,\, 0.0]$ & $\log_{10} f_{\mathrm{esc},10}$ (escape norm.) \\
\texttt{ALPHA\_ESC} & $[-1.0,\, 0.5]$ & $\alpha_{\mathrm{esc}}$ (escape slope) \\
\texttt{M\_TURN} & $[8.0,\, 11.0]$ & $\log_{10}(M_{\mathrm{turn}}/M_\odot)$ \\
\texttt{t\_STAR} & $[0.01,\, 1.0]$ & $t_\star$ (SF timescale) \\
\texttt{L\_X} & $[38.0,\, 42.0]$ & $\log_{10}(L_X/\mathrm{SFR})$ ($\mathrm{erg\,s^{-1}\,yr}\,M_\odot^{-1}$) \\
\texttt{NU\_X\_THRESH} & $[100,\, 1500]$ & $E_0/\mathrm{eV}$ (X-ray threshold) \\
\texttt{KP\_h\_Mpc} & $[1.0,\, 30.0]$ & $k_p/(h\,\mathrm{Mpc}^{-1})$ (break) \\
\texttt{MS} & $[0.5,\, 4.0]$ & $m_s$ (small-scale slope) \\
\end{tabular}
\end{ruledtabular}
\end{table}

\subsection{Simulation Contract, Observables, and GSL Thread-Count Reproducibility}\label{sec:contract}

All production simulations use a cubic comoving box of length $\texttt{BOX\_LEN}=250\,\mathrm{Mpc}$ with $\texttt{DIM}=512$ high-resolution initial density/velocity grid cells and $\texttt{HII\_DIM}=128$ ionization/brightness-temperature cells, second-order Lagrangian perturbation theory (\texttt{USE\_2LPT=true}), interpolation tables (\texttt{USE\_INTERPOLATION\_TABLES=true}), inhomogeneous recombinations (\texttt{INHOMO\_RECO=true}), mass-dependent ionizing efficiency (\texttt{USE\_MASS\_DEPENDENT\_ZETA=true}), spin-temperature fluctuations (\texttt{USE\_TS\_FLUCT=true}), redshift-space distortions (\texttt{APPLY\_RSDS=true}), and Ly$\alpha$/CMB heating enabled, compiled in a frozen native library binary (SHA-256 prefix \texttt{3211a6629109da73}). Coeval outputs are extracted on a standardized 32-node redshift grid $z \in \{5.0, 5.9, 6, 7, \dots, 35\}$, recording both volume-averaged neutral fraction $x_{\mathrm{HI}}(z)$ and spherically binned 21-cm brightness-temperature dimensionless power spectra $\Delta^2_{21}(k,z)$ in $\mathrm{mK}^2$.

Under the joint \texttt{ps\_xhi\_fullbox\_v1} data contract, full-box spherical FFT power spectra are first binned onto 60 uniform shells $\texttt{ps\_native}$ up to the $128^3$ corner Nyquist mode ($k \in [0.0232, 2.7628]\,\mathrm{Mpc}^{-1}$) and resampled without extrapolation via piecewise-linear interpolation in $k$ onto 32 logarithmically spaced evaluation centers $\texttt{ps\_full}$ spanning $k \in [0.1, 1.0]\,\mathrm{Mpc}^{-1}$ ($\mathrm{mK}^2$). Each design point stores paired coeval $(x_{\mathrm{HI}}, \Delta^2_{21})$ from identical density and ionization boxes in atomic per-point \texttt{NPZ} archives carrying design index, requested and effective random IC seeds, split metadata, and configuration SHA-256 hashes.

A critical numerical reproducibility finding discovered during cross-validation of CPU and GPU backends impacts all downstream emulator dataset audits: native \textsc{21cmFAST} partitions GSL pseudo-random number generator streams across OpenMP threads during initial Fourier-mode generation. Consequently, \emph{identical integer \texttt{random\_seed} values generate different initial Fourier density realizations when executed with $\texttt{N\_THREADS}=8$ versus $\texttt{N\_THREADS}=16$}. In paired CPU--GPU cross-checks weighted by SKA-Low diagonal uncertainties, an 8-thread replay against 16-thread history labels produced spurious $3\text{--}13\sigma$ single-bin residuals (similarly present in 8-thread pure-CPU reruns), whereas matching $\texttt{N\_THREADS}=16$ lowered maximum single-bin backend residuals to $0.033\text{--}0.038\sigma$ across four-IC test realizations and $0.052\sigma$ on fixed-seed comparisons. All canonical emulator data contracts, posterior-stratified direct-simulation audits (Sec.~\ref{sec:a3}), and full-box production archives therefore explicitly pin $\texttt{N\_THREADS}=16$ alongside integer seed and native binary SHA-256.

\subsection{Observational Likelihood Interface}\label{sec:likelihood}

Inference interfaces combine an exact analytical/numerical UV LF likelihood $\ell_{\mathrm{LF}}^{\mathrm{exact}}(\theta)$ (spanning 34 HST \cite{Bouwens2015,Bouwens2017} and JWST \textsc{PRIMER} \cite{Donnan2024} points at $z=6\text{--}10$ with $M_{\mathrm{UV}} > -20$ without emulator approximation) with two reionization observables computed directly from the emulated 32-node neutral-fraction history $\widehat{X}(\theta) \equiv \{\widehat{x}_{\mathrm{HI}}(z_j)\}_{j=1}^{32}$:
\begin{align}
\ell_{\tau}[X] &= -\frac{1}{2}\left[\frac{\tau[X] - 0.0544}{\sigma_{\tau}(X)}\right]^2,\nonumber\\
\sigma_\tau(X) &= \begin{cases} 0.0081, & \tau[X] < 0.0544,\\ 0.0070, & \tau[X] \ge 0.0544, \end{cases}\\
\ell_x[X] &= -\frac{1}{2}\left[\frac{\max\{0,\, X(5.9) - 0.06\}}{0.05}\right]^2,
\end{align}
corresponding respectively to the \textit{Planck} 2018 TT,TE,EE+lowE Thomson scattering optical depth constraint \cite{Planck2018} and the McGreer \textit{et al.} dark-pixel neutral-fraction limit at $z=5.9$ \cite{McGreer2015}. The Thomson optical depth functional $\tau[X]$ is integrated using the executable \texttt{compute\_tau} adapter (\texttt{tau\_from\_history}) over 31 uniformly spaced redshift nodes spanning $z \in [5, 35]$ interpolated from the 32-node network output vector, while the additional $z=5.9$ node retained in that 32-node vector is evaluated directly for $\ell_x[X]$. Defining $\ell_{\mathrm{hist}}[X] \equiv \ell_{\tau}[X] + \ell_x[X]$, the deployed representative-history log-likelihood reads
\begin{equation}
\ell_{\mathrm{rep}}(\theta) = \ell_{\mathrm{LF}}^{\mathrm{exact}}(\theta) + \ell_{\mathrm{hist}}[\widehat{X}(\theta)].
\label{eq:lrep}
\end{equation}

\section{Ten-Dimensional Neutral-Fraction History Emulator}\label{sec:history}

\subsection{Architecture and Ensembling}\label{sec:xhi_arch}

The history-emulator architecture is governed by the frozen \texttt{architecture\_contract.json} specification: a pre-normalized residual multilayer perceptron (PreNorm ResMLP) mapping ten z-scored physical coordinates to the 32-node vector $x_{\mathrm{HI}}(z) \in (0, 1)^{32}$ via a final sigmoid decoder, without heuristic monotonicity or end-of-reionization clamping so that incomplete reionization histories produced natively by the simulator at $z=5$ remain faithfully represented. The backbone comprises an input linear projection of width $W=256$, four residual blocks each applying LayerNorm, two linear transforms of hidden width $H=512$ with $\mathrm{SiLU}$ activation, and residual contribution scaled by $0.5$, followed by final LayerNorm and linear output projection---totaling $1{,}065{,}248$ trainable parameters per network ($5{,}326{,}240$ across the five-seed ensemble)---initialized with Xavier uniform weights and zero biases. Standardization statistics are computed strictly from the training split. Networks minimize physical-space mean squared error (MSE) equally weighted across all 32 redshift nodes using \texttt{AdamW} (learning rate $10^{-3}$, weight decay $10^{-4}$, batch size $128$, up to $1{,}000$ epochs with validation-plateau learning-rate decay and early stopping). To suppress initialization variance and smooth predictions, five independent random training seeds ($11, 29, 47, 71, 101$) are averaged arithmetically in physical $x_{\mathrm{HI}}$ space. In runtime logs, a single five-seed surrogate prediction evaluates in $\sim 0.41\,\mathrm{ms}$ ($407.3\,\mu\mathrm{s}$ at batch size 1 on 4 CPU threads) compared with a median $416\,\mathrm{s}$ ($q_{90}=494\,\mathrm{s}$) wall time per 16-thread direct \textsc{21cmFAST} box, while full Bayesian inference chains incorporate exact UV LF evaluation and optical-depth postprocessing alongside history prediction. Each training design point in the scrambled Sobol design (\texttt{random\_ic\_v2}) carries a single independent random IC seed that is withheld from network inputs, causing MSE minimization to target the conditional expectation $\mathbb{E}_{\mathrm{IC}}[x_{\mathrm{HI}} \mid \theta]$.

\subsection{Accuracy Verification and 65,536-to-100,000 Scaling}\label{sec:xhi_acc}

Figure~\ref{fig:xhi_acc} and Table~\ref{tab:xhi_val} quantify trajectory and observable reconstruction across two model generations: the initial 65,536-history baseline (\texttt{20261026}) and the expanded 100,000-history production surrogate (\texttt{random\_ic\_development\_100000\_20261103}). Because the frozen 512-history benchmark and full 10,000-history split serve as development-validation splits used to monitor learning-rate schedules and early stopping, they measure surrogate convergence under matched prior sampling, whereas strictly disjoint out-of-sample generalization is established by the 56 posterior direct simulations in Sec.~\ref{sec:a3} (with a further 5,000 sealed design points retained unread and unevaluated under contract isolation).

On the paired 512-history development benchmark (exact matched parameter IDs and reference labels verified via SHA-256), averaging five seeds outperforms every individual member: individual seeds span median trajectory RMSE $0.001276\text{--}0.001524$ ($q_{95} = 0.007786\text{--}0.010080$), whereas the five-seed ensemble lowers median RMSE to $0.001008$ and $q_{95}$ to $0.006815$ (mean $0.002132$, maximum $0.040523$; $q_{95}$ within-history max absolute error $0.028465$). Scaling training coverage by $52.6\%$ from 65,536 to 100,000 unique parameter-IC pairs on those identical 512 validation points reduces every downstream error metric monotonically: $x_{\mathrm{HI}}(5.9)$ RMSE drops by $35.6\%$ ($0.006986 \to 0.004501$, MAE $-22.7\%$), $\tau$ RMSE drops by $23.2\%$ ($8.74\times 10^{-4} \to 6.72\times 10^{-4}$, $q_{95}$ $-27.9\%$), and reionization log-likelihood errors $|\Delta\ell_{\mathrm{hist}}|$ drop by $26.6\%$ at $q_{90}$ ($1.612 \to 1.183$) and $25.2\%$ at $q_{95}$ ($2.734 \to 2.046$). Evaluated over all 10,000 development-validation points of the 100k model, median trajectory RMSE reaches $8.51\times 10^{-4}$ ($q_{95} = 5.56\times 10^{-3}$), $q_{90}$ error in $\tau$ is $3.35\times 10^{-4}$, $q_{90}$ error in $x_{\mathrm{HI}}(5.9)$ is $0.00457$ with only $25/10{,}000$ threshold crossings across $x_{\mathrm{HI}}(5.9)=0.06$, and prior-box $q_{90}(|\Delta\ell_{\mathrm{hist}}|)$ is $1.177$.

\figwide{figures/fig1_history_accuracy.pdf}{Validation of the 10D history emulator across the frozen 512-history development benchmark and 100,000-history expansion. (a) Representative neutral-fraction trajectories at ascending-RMSE ranks 256 (near-median, $\mathrm{RMSE}=0.00100$; full 512-point median $0.00101$) and 487 ($q_{95}$, $\mathrm{RMSE}=0.00690$), comparing simulated truth (solid) with five-seed ensemble predictions (dashed) and five-seed envelope (shading) over $z \in [5,22]$. (b) Signed residuals $(\widehat{x}_{\mathrm{HI}}-x_{\mathrm{HI}})\times 10^3$ compared against the 512-history $q_{90}$ envelope (blue band) and the 128-family prior-box IC sensitivity audit (\texttt{T7}) mean IC scatter $\bar{s}_{\mathrm{IC}}(z)$ (dot-dashed). (c) Empirical trajectory-RMSE cumulative distribution across all 512 histories showing the variance reduction from five individual seeds (thin gray) to their ensemble (thick navy), with 100k 10,000-history development quantiles indicated vertically. (d) Paired percentage error reductions on identical 512 development-validation histories when scaling from 65,536 to 100,000 training realizations.}{fig:xhi_acc}

\subsection{Finite Exact-vs.-Emulator Inference Smoke Check}\label{sec:t6}

As an early finite-sample exact-vs.-emulator Markov-chain check (denoted \texttt{T6}), we compared a 32-walker, 200-step direct-\textsc{21cmFAST} smoke chain (steps 100--199, 3,200 correlated rows at fixed $(k_{p,h}, m_s)=(10\,h\,\mathrm{Mpc}^{-1}, 2.5)$, $L_X=40.5$, $E_0=500\,\mathrm{eV}$) against steps 2500--4999 of a 5,000-step 65k-emulator chain (80,000 rows). Across the six sampled galaxy parameters ($\eta_\star, t_\star, \alpha_\star, \log_{10} M_{\mathrm{turn}}, \log_{10} f_{\mathrm{esc},10}, \alpha_{\mathrm{esc}}$), standardized finite-sample mean offsets $(\mu_e-\mu_m)/\sigma_m$ remain within $[-0.100, +0.219]$ ($\alpha_{\mathrm{esc}}$ largest) and standard-deviation ratios $s_e/s_m$ span $0.895\text{--}1.068$ with correlation-matrix Frobenius difference $0.669$ (Fig.~\ref{fig:t6}). Because the 200-step exact segment is explicitly tagged unconverged and non-equilibrium (\texttt{SMOKE\_ONLY}, establishing neither target parity nor asymptotic posterior convergence), \texttt{T6} provides only descriptive finite-sample consistency and motivates the direct posterior-stratified out-of-sample \textsc{21cmFAST} validations below.

\figcolumn{figures/fig2_t6_exact_vs_emu.pdf}{Descriptive finite-chain exact-vs.-emulator smoke comparison (\texttt{T6}) between a 3,200-row direct-\textsc{21cmFAST} run and an 80,000-row 65,536-model emulator segment across six astrophysical parameters at fixed $(k_{p,h}, m_s)=(10\,h\,\mathrm{Mpc}^{-1}, 2.5)$. (a) Mean difference normalized by the exact finite-sample standard deviation ($\max = 0.219$). (b) Sample width ratios ($0.895\text{--}1.068$; shaded bands mark $\pm 0.25\sigma$ and $\pm 10\%$).}{fig:t6}

\begin{table*}[!tb]
\caption{Validation performance of the history emulator. Top: five-seed vs.\ individual-seed trajectory RMSE on the frozen 512-history split (65,536 model). Middle: paired comparison of 65,536 vs.\ 100,000 models on the same 512 validation points. Bottom: full 10,000-history development summary of the 100k model. Superscript $a$ denotes $q_{95}$ (95th percentile); superscript $b$ denotes the maximum trajectory RMSE across the 512 points (mean is $0.002132$).}
\label{tab:xhi_val}
\footnotesize
\begin{ruledtabular}
\begin{tabular}{llccc}
Model / split & Metric & Median / MAE & $q_{90}$ / $q_{95}$ & RMSE / max \\
\hline
65k, 5-seed ensemble (512 val) & Trajectory RMSE & 0.001008 & $0.006815^{a}$ & $0.04052^{b}$ \\
65k, individual seeds (512 val) & Trajectory RMSE & 0.001276--0.001524 & 0.007786--$0.010080^{a}$ & \dots \\
\hline
65,536 paired (512 val) & $x_{\mathrm{HI}}(5.9)$ & 0.002289 & $0.005206 / 0.009286$ & 0.006986 \\
100,000 paired (512 val) & $x_{\mathrm{HI}}(5.9)$ & 0.001769 & $0.004427 / 0.008188$ & 0.004501 \\
65,536 paired (512 val) & $\tau\ (\times 10^4)$ & 2.637 & $4.415 / 9.190$ & 8.742 \\
100,000 paired (512 val) & $\tau\ (\times 10^4)$ & 2.106 & $3.684 / 6.627$ & 6.717 \\
65,536 paired (512 val) & $|\Delta\ell_{\mathrm{hist}}|$ & 0.896 & $1.612 / 2.734$ & 3.494 \\
100,000 paired (512 val) & $|\Delta\ell_{\mathrm{hist}}|$ & 0.728 & $1.183 / 2.046$ & 3.145 \\
\hline
100k dev split (10,000) & Trajectory RMSE & 0.000851 & $0.005564^{a}$ & \dots \\
100k dev split (10,000) & $x_{\mathrm{HI}}(5.9)$ & 0.001888 & $0.004571 / 0.008630$ & 0.005698 \\
100k dev split (10,000) & $\tau\ (\times 10^4)$ & \dots & 3.351 ($q_{90}$) & \dots \\
100k dev split (10,000) & $|\Delta\ell_{\mathrm{hist}}|$ & \dots & 1.177 ($q_{90}$) & \dots \\
\end{tabular}
\end{ruledtabular}
\end{table*}

\section{Random Initial Conditions and Likelihood Target Semantics}\label{sec:ic}

Because the history emulator trains on single independent random IC realizations without conditioning on the IC seed, its regression target approximates the conditional mean $\bar{X}(\theta) = \mathbb{E}_{\mathrm{IC}}[x_{\mathrm{HI}} \mid \theta]$. Because the likelihood $\mathcal{L}_{\mathrm{hist}}[X] \equiv \exp(\ell_{\mathrm{hist}}[X])$ is nonlinear---especially near the piecewise-quadratic threshold $x_{\mathrm{HI}}(5.9)=0.06$ and in exponentially steep prior-box tails---Jensen's inequality and realization variance imply
\begin{equation}
\mathbb{E}_{\mathrm{IC}}\!\left[\mathcal{L}_{\mathrm{hist}}(X)\right]
\neq \mathcal{L}_{\mathrm{hist}}\!\left(\mathbb{E}_{\mathrm{IC}}[X]\right).
\label{eq:jensen}
\end{equation}
Evaluating emulator fidelity therefore requires distinguishing prior-volume worst-case scatter from likelihood accuracy inside the region supported by observations.

\subsection{Prior-Box Multi-IC Sensitivity Audit (T7)}\label{sec:t7}

We evaluated 128 parameter families across the full ten-dimensional design box with eight realizations per family ($128$ fixed references plus $896$ fresh random ICs; denoted \texttt{T7} in Table~\ref{tab:t7_a3}). Define the family eight-IC average history $\bar{X}_8$ and the signed gap between representative-history and IC-marginal log-likelihoods,
\begin{equation}
\delta\ell \equiv \ell_{\mathrm{hist}}[\bar{X}_8] - \ln\!\left[\frac{1}{8}\sum_{r=1}^8 \exp\left(\ell_{\mathrm{hist}}[X_r]\right)\right].
\label{eq:deltal}
\end{equation}
Across the full prior box, physical trajectory IC scatter is modest ($q_{50}=2.88\times 10^{-4}$, $q_{90}=0.00313$, $s_{\mathrm{IC}}(\tau)$ $q_{90}=1.35\times 10^{-4}$, $s_{\mathrm{IC}}[x_{\mathrm{HI}}(5.9)]$ $q_{90}=0.00751$, with mean IC scatter peaking at $0.00213$ at $z=5.9$). However, in low-$f_{\mathrm{esc}}$ or late-reionizing prior draws where $\ell_{\mathrm{hist}} \ll -10$, small absolute variations in $x_{\mathrm{HI}}(5.9)$ amplify into large log-likelihood spread: $q_{90}[s_{\mathrm{IC}}(\ell_{\mathrm{hist}})] = 0.948$ ($\max = 2.210$) and $13/128$ families cross the $x_{\mathrm{HI}}(5.9)=0.31$ diagnostic screen across realizations, failing the conservative uniform-prior single-IC screen (\texttt{MULTI\_IC\_AVERAGING\_REQUIRED}). Furthermore, reporting signed quantiles of $\delta\ell$ ($q_{50}=-0.00249$, $q_{90}=1.12\times 10^{-7}$, $q_{95}=6.03\times 10^{-7}$) understates tail magnitude because large positive and negative discrepancies cancel around zero; the completion audit records $q_{90}(|\delta\ell|) = 0.308$ and $\max |\delta\ell| = 1.908$.

\begin{table*}[!tb]
\caption{Prior-box eight-IC sensitivity audit (\texttt{T7}: 128 families, 1,024 realizations) contrasted with the out-of-sample posterior-stratified direct-\textsc{21cmFAST} audit (\texttt{A3}: 32 unique points, 56 tasks, $\texttt{N\_THREADS=16}$). Note $a$: reported in the \texttt{T7} completion audit; the released summary JSON lists signed quantiles $q_{50}=-0.002491$, $q_{90}=1.12\times 10^{-7}$, $q_{95}=6.03\times 10^{-7}$.}
\label{tab:t7_a3}
\begin{ruledtabular}
\begin{tabular}{llccc}
Audit / subset & Quantity & Median / MAE & $q_{90}$ & Max \\
\hline
\texttt{T7} prior (128 families) & Trajectory $s_{\mathrm{IC}}$ RMS & 0.000288 & 0.003133 & 0.005158 \\
\texttt{T7} prior (128 families) & $s_{\mathrm{IC}}(\tau)\ (\times 10^4)$ & 0.286 & 1.351 & 2.736 \\
\texttt{T7} prior (128 families) & $s_{\mathrm{IC}}[x_{\mathrm{HI}}(5.9)]$ & 0.000125 & 0.007511 & 0.02205 \\
\texttt{T7} prior (128 families) & $s_{\mathrm{IC}}(\ell_{\mathrm{hist}})$ & 0.0769 & 0.9481 & 2.210 \\
\texttt{T7} prior (128 families) & $|\delta\ell|$ (mean history vs.\ IC avg.) & \dots & $0.3082^{a}$ & 1.908 \\
\hline
\texttt{A3} \texttt{joint8d} (28 tasks) & $|\Delta\ell_{\mathrm{reion}}|$ & 0.024 & 0.056 & 0.153 \\
\texttt{A3} \texttt{fixed\_joint} (20 tasks) & $|\Delta\ell_{\mathrm{reion}}|$ & 0.011 & 0.026 & 0.042 \\
\texttt{A3} \texttt{fixed\_lf\_tau} (8 tasks) & $|\Delta\ell_{\tau}|$ (active $\tau$ term) & 0.010 & 0.029 & 0.036 \\
\texttt{A3} \texttt{fixed\_lf\_tau} (8 tasks) & $|\Delta\ell_{\mathrm{reion}}|$ (unused $x_{\mathrm{HI}}$ tail) & 0.469 & 1.315 & 1.336 \\
\texttt{A3} 4-IC posterior (8 pts) & $|\ell(\bar{X}_4)-\ell_{\mathrm{IC\text{-}avg}}|$ & \dots & \dots & 0.0054 \\
\texttt{A3} 4-IC posterior (8 pts) & $|\ell_{\mathrm{emu}}-\ell_{\mathrm{IC\text{-}avg}}|$ & \dots & \dots & 0.0274 \\
\texttt{A3} 4-IC posterior (8 pts) & $|\ell_{\mathrm{emu}}-\ell(\bar{X}_4)|$ & \dots & \dots & 0.0328 \\
\end{tabular}
\end{ruledtabular}
\end{table*}

\subsection{Reconciliation Inside Observational Posterior Support (A3)}\label{sec:a3}

Does prior-box IC amplification distort Bayesian posteriors dominated by UV LF + $\tau$ + $x_{\mathrm{HI}}(5.9)$ observations? To test out-of-sample fidelity directly, we executed the stratified posterior audit (\texttt{A3}) at 32 unique parameter points drawn from converged 8D (\texttt{joint8d}) and 6D (\texttt{fixed\_joint}, \texttt{fixed\_lf\_tau}) 100k history-emulator posterior chains---explicitly oversampling bulk high-probability regions, low-$f_{\mathrm{esc}}$ tails, and the $x_{\mathrm{HI}}(5.9)=0.06$ kink---totaling 56 direct \textsc{21cmFAST} simulations pinned to \texttt{N\_THREADS=16} (including eight points simulated with four independent ICs each; Table~\ref{tab:t7_a3} and Fig.~\ref{fig:ic_post}).

Inside actual posterior support, likelihood errors collapse by more than an order of magnitude relative to the prior box: $\ell_{\mathrm{reion}}$ achieves $\mathrm{MAE}=0.024$ and $q_{90}=0.056$ ($\max = 0.153$) across 28 \texttt{joint8d} tasks, and $\mathrm{MAE}=0.011$, $q_{90}=0.026$ ($\max = 0.042$) across 20 \texttt{fixed\_joint} tasks. By posterior stratum, the two large $|\Delta\ell_{\mathrm{reion}}|$ excursions occur exclusively in \texttt{fixed\_lf\_tau} runs where the $x_{\mathrm{HI}}(5.9)$ penalty is turned off, reaching $1.305$ in the \texttt{bulk} stratum ($q_{90}=0.026$ across all 28 \texttt{bulk} tasks) and $1.336$ in \texttt{low\_fesc\_tail}, whereas the active $\ell_{\tau}$ term in \texttt{fixed\_lf\_tau} remains accurate within $\max = 0.036$. Most decisively, across all eight four-IC posterior points---including four points straddling the $x_{\mathrm{HI}}(5.9)=0.06$ threshold---individual IC realizations span $[-0.055, +0.035]$ ($\max |\ell - \ell_{\mathrm{IC\text{-}avg}}| = 0.055$) around $\ell_{\mathrm{IC\text{-}avg}}$, whereas replacing IC marginalization by the four-IC mean history $\ell(\bar{X}_4)$ differs by at most $0.0054$, and the 100k history-emulator prediction tracks $\ell_{\mathrm{IC\text{-}avg}}$ within $\le 0.0274$ ($\le 0.028$ in Fig.~\ref{fig:ic_post}b) and tracks $\ell(\bar{X}_4)$ within $\le 0.0328$ ($\le 0.033$; Table~\ref{tab:t7_a3}). This directly confirms that representative-history emulation introduces negligible bias within the observational posterior support while avoiding multi-IC likelihood averaging at every MCMC step.

\figwide{figures/fig3_ic_and_posterior_val.pdf}{Reconciling prior-box \texttt{T7} IC variance with out-of-sample \texttt{A3} posterior direct-\textsc{21cmFAST} validation. (a) Magnitude of log-likelihood discrepancies ($|\Delta \ell|$ on log scale) across three prior-box diagnostics (shaded upper rows: \texttt{T7} IC spread, \texttt{T7} $|\delta\ell|$ mean-history vs.\ IC-average with $q_{90}$ from completion audit, and 100k 512-history validation) compared against \texttt{A3} stratified posterior runs (\texttt{joint8d}, \texttt{fixed\_joint}, \texttt{fixed\_lf\_tau} active $\ell_\tau$, and 4-IC $|\ell(\bar{X}_4)-\ell_{\mathrm{avg}}|$). The cross marks the inactive low-$f_{\mathrm{esc}}$ $x_{\mathrm{HI}}$ tail not used in \texttt{fixed\_lf\_tau}. (b) Point-by-point offsets across all eight four-IC posterior points comparing individual IC scatter (slate dots/bars, spanning $[-0.055, +0.035]$), $\ell(\bar{X}_4)$ (navy squares, $\le 0.0054$), and 100k emulator predictions (amber diamonds, $\le 0.028$ vs.\ $\ell_{\mathrm{IC\text{-}avg}}$).}{fig:ic_post}

\section{Data-Engineering Design for Full-Box 21-cm Power-Spectrum Archives}\label{sec:ps}

Extending surrogate modeling from scalar neutral-fraction curves to the 21-cm brightness-temperature power spectrum $\Delta^2_{21}(k,z)$ over 32 redshifts and 32 $k$-bins requires diagnosing two fundamental data-engineering trade-offs before generating full-scale training suites: (i) how to allocate finite simulation budgets between unique design points and IC replicates per point, and (ii) whether subbox power-spectrum averages can substitute for full-box spherical FFTs. Throughout this section, power-spectrum errors are evaluated as root-mean-square error in $\log_{10}(\Delta^2_{21}/\mathrm{mK}^2 + 0.01)$ (denoted in $\mathrm{dex}$), preventing divide-by-zero blowup when reionization nears completion.

\subsection{Single-IC Coverage vs.\ Multi-IC Replication}\label{sec:ps_ic_budget}

Using an initial $9\times 20$ redshift-wavenumber pilot ($z \in [6.34, 18.21]$, $k \in [0.116, 0.996]\,\mathrm{Mpc}^{-1}$, 180 $(z,k)$ cells, with $N_{\mathrm{val}}=103$ and held-out four-IC mean spectra across $N_{\mathrm{test}}=108$ independent points), we isolated the impact of training with one IC versus four ICs per parameter point at fixed $N_{\mathrm{train}}=534$, as well as scaling $N_{\mathrm{train}}$ at fixed architecture (Table~\ref{tab:ps_pilot} and Fig.~\ref{fig:ps_scaling}). Critically, the subbox pilot archive samples only a restricted small-scale slope sub-domain $m_s \in [1.51, 4.0]$ (exact contract support $[1.5076, 4.0]$), omitting the flatter/redder branch $m_s \in [0.5, 1.51]$ encompassing the $\Lambda$CDM limit ($m_s=0.968$). Within that shared $[1.51, 4.0]$ test domain, pilot comparisons isolate sampling and IC allocation cleanly.

At $N_{\mathrm{train}}=534$, single-IC training achieves $0.5808\,\mathrm{dex}$ overall log-RMSE (median fractional error $41.1\%$, $q_{90}=193.7\%$), nearly indistinguishable from four-IC target averaging ($0.5786\,\mathrm{dex}$, median $41.0\%$, $q_{90}=192.9\%$) despite costing one quarter as many simulations. Across the $(z,k)$ plane, emulator RMSE exceeds four-IC intrinsic sample scatter $s_{\mathrm{IC}}$ by factors of $4.89\text{--}65.8\times$ (Fig.~\ref{fig:ps_scaling}a), peaking near mid-reionization $z \sim 8\text{--}10$. Conversely, expanding parameter coverage from $267$ (4-IC) to $534$ (4-IC) to $2,326$ (mixed 1-IC/4-IC) lowers overall RMSE sharply from $0.5965$ to $0.5786$ to $0.3435\,\mathrm{dex}$ (and low-$z$ RMSE from $0.7434 \to 0.7284 \to 0.4450\,\mathrm{dex}$; Fig.~\ref{fig:ps_scaling}b). Broad 10D parameter coverage with one random IC per point thus strictly dominates multi-IC replication under any fixed simulation budget.

\figwide{figures/fig4_ps_pilot_scaling_and_ic.pdf}{Controlled pilot findings guiding the full-box 21-cm power-spectrum data strategy. (a) Log-RMSE ($\log_{10}[\Delta^2_{21}+0.01]$ in dex, color shading) overlaid with contours of the emulator-to-4-IC-scatter ratio ($\mathrm{RMSE}/s_{\mathrm{IC}}$, spanning $4.9\times\text{--}65.8\times$) over the $9\times 20$ pilot grid ($m_s \in [1.51, 4.0]$) at $N_{\mathrm{train}}=534$. (b) Overall (navy) and low-$z$ ($z\le 7.97$, dashed teal) test RMSE vs.\ unique training design points $N_{\mathrm{train}}$, demonstrating parity between 1-IC ($0.581\,\mathrm{dex}$, amber diamond) and 4-IC ($0.579\,\mathrm{dex}$) training at $N=534$ alongside monotonic coverage gains to $N=2{,}326$ ($0.344\,\mathrm{dex}$), contrasted with the $N=128$ full-domain ($m_s \in [0.5, 4.0]$) $32\times 32$ scratch MLP ($1.129\,\mathrm{dex}$, amber square) and constant baseline ($1.437\,\mathrm{dex}$, cross).}{fig:ps_scaling}

\begin{table*}[!tb]
\caption{Power-spectrum pilot scaling and estimator-bias experiments evaluated in $\log_{10}(\Delta^2_{21}+0.01)$ ($\mathrm{dex}$), explicitly separating restricted-slope ($m_s \in [1.51, 4.0]$) subbox pilot tests from full-domain ($m_s \in [0.5, 4.0]$) 250\,Mpc full-box evaluations.}
\label{tab:ps_pilot}
\footnotesize
\setlength{\tabcolsep}{3.5pt}
\begin{ruledtabular}
\begin{tabular}{llccc}
Experiment / baseline & Grid ($z \times k$) & Train / test $m_s$ & Held-out test set & Overall RMSE (dex) \\
\hline
Pilot $N=267$ (4-IC training) & $9 \times 20$ (subbox) & $[1.51,4.0] \to [1.51,4.0]$ & 108 (4-IC mean) & 0.5965 \\
Pilot $N=534$ (1-IC training) & $9 \times 20$ (subbox) & $[1.51,4.0] \to [1.51,4.0]$ & 108 (4-IC mean) & 0.5808 \\
Pilot $N=534$ (4-IC training) & $9 \times 20$ (subbox) & $[1.51,4.0] \to [1.51,4.0]$ & 108 (4-IC mean) & 0.5786 \\
Pilot $N=2{,}326$ (mixed 1-IC / 4-IC) & $9 \times 20$ (subbox) & $[1.51,4.0] \to [1.51,4.0]$ & 108 (4-IC mean) & 0.3435 \\
Subbox model vs.\ full-box truth & $9 \times 20$ & $[1.51,4.0] \to [0.5,4.0]$ & 64 (full-box 4-IC) & 1.0661 \\
$\hookrightarrow$ after 32-pair median-ratio corr. & $9 \times 20$ & $[1.51,4.0] \to [0.5,4.0]$ & 64 (full-box 4-IC) & 1.0631 \\
Full-box constant baseline & $32 \times 32$ (full) & $[0.5,4.0]$ & 64 (full-box 4-IC) & 1.4366 \\
Full-box scratch MLP ($N=128$) & $32 \times 32$ (full) & $[0.5,4.0] \to [0.5,4.0]$ & 64 (full-box 4-IC) & 1.1286 \\
\end{tabular}
\end{ruledtabular}
\end{table*}

\subsection{Subbox Estimator Bias and Full-Box Contract Production}\label{sec:ps_estimator}

A second architectural hazard arises from evaluating power spectra via 12-subbox averaging rather than the native 250\,Mpc full-box spherical FFT. Direct paired comparisons at $k \in [0.1, 1.0]\,\mathrm{Mpc}^{-1}$ expose strongly parameter-dependent low-$k$ distortion (Fig.~\ref{fig:ps_est}a): while design point \texttt{199} ($k_{p,h}=4.33$, $m_s=4.00$) stays within $0.95\text{--}1.05\times$, point \texttt{833} ($k_{p,h}=1.62$, $m_s=3.27$) exhibits a monotonic excess reaching $1.480\times$ at $z=7.97, k=0.116\,\mathrm{Mpc}^{-1}$ ($1.508\times$ at $k=0.101\,\mathrm{Mpc}^{-1}$). Evaluating the 2,326-point subbox-trained emulator ($m_s \in [1.51, 4.0]$) against 64 independent four-IC 250\,Mpc full-box spectra spanning the broader full contract ($m_s \in [0.5, 4.0]$) yields $1.0661\,\mathrm{dex}$ overall error ($0.7525\,\mathrm{dex}$ at low $z$ and $0.8654\,\mathrm{dex}$ at low $z$, low $k$). That $0.3435 \to 1.0661\,\mathrm{dex}$ jump deliberately reflects two superposed failures: estimator mismatch between subbox and full-box FFTs plus out-of-domain extrapolation into $m_s \in [0.5, 1.51]$. Even with a 32-pair empirical median-ratio correction, full-box error shifts only marginally to $1.0631\,\mathrm{dex}$ (with low-$z$, low-$k$ error moving slightly from $0.8654$ to $0.8671\,\mathrm{dex}$). Matching training and evaluation domains ($m_s \in [0.5, 4.0]$) and estimators natively via full-box spherical FFTs is therefore mandatory.

Training on the full $32\times 32$ contract grid with only $N=128$ initial full-box realizations improves over the constant baseline ($1.4366 \to 1.1286\,\mathrm{dex}$) while visibly underfitting spectral shapes at median-error validation point \texttt{80475} (Fig.~\ref{fig:ps_est}b). We therefore executed full-scale simulation production under the frozen \texttt{ps\_xhi\_fullbox\_v1} data contract pinned to \texttt{N\_THREADS=16}: all $100{,}000 / 100{,}000$ single-IC 250\,Mpc full-box spherical-FFT realizations have completed generation (with a frozen $95{,}000$ train / $5{,}000$ development export manifest containing 14 verified replacements) alongside 64 independent four-IC full-box validation realizations, establishing a verified full-box $32 \times 32$ archive to support future power-spectrum emulator development and release auditing.

\figwide{figures/fig5_ps_estimator_and_fullbox.pdf}{Subbox estimator bias and full-box $32\times 32$ validation baseline. (a) Ratio of 12-subbox to 250\,Mpc full-box spherical-FFT power spectra across $k \in [0.1, 1.0]\,\mathrm{Mpc}^{-1}$ for design point \texttt{833} ($z=6.34, 7.10, 7.97$) and design point \texttt{199} ($z=6.34, 7.97$), highlighting the strongly parameter-dependent low-$k$ excess (up to $1.48\times$ at $k=0.116\,\mathrm{Mpc}^{-1}$) that prevents global ratio rescaling of subbox models. (b) Representative $32\times 32$ full-box power spectra ($z=7, 8, 10$) for the median-error validation box (\texttt{id=80475}) comparing true four-IC full-box spectra (solid) against the under-sampled $N=128$ initial baseline (dashed, $1.129\,\mathrm{dex}$ overall), illustrating the target dynamic range addressed by our 100,000-box production dataset.}{fig:ps_est}

\section{Joint Power-Spectrum and History Emulation (\texorpdfstring{\textsc{PSX-emu v1.0}}{PSX-emu v1.0}) and Localized Residual Patching}\label{sec:psx_emu}

\subsection{Multi-Task Joint Architecture, Split-Contamination Audit, and 300-Point Independent Evaluation}\label{sec:psx_audit}

Using the completed 100,000-realization full-box \texttt{ps\_xhi\_fullbox\_v1} archive ($32 \times 32$ power spectra spanning $z \in \{5.0, 5.9, 6, \dots, 35\}$ and $k \in [0.1, 1.0]\,\mathrm{Mpc}^{-1}$ paired with 32-node $x_{\mathrm{HI}}(z)$), we train a shared-trunk multi-task residual emulator (\textsc{PSX-emu v1.0}) that simultaneously predicts $\log_{10}(\Delta^2_{21}/\mathrm{mK}^2 + 0.01)$ and $x_{\mathrm{HI}}(z)$ across all ten parameters ($\mathrm{base}=128$, $878{,}209$ trainable parameters per network, compared against $1{,}065{,}248$ in the standalone history-only network). We tested six loss-weighting configurations using three-seed ensembles trained for $800$ epochs on the historical 64-point validation split: fixed joint weighting $\lambda=1.0$ (\texttt{emu\_joint\_l10}) achieved the best joint balance ($\mathrm{RMSE}_{\log_{10}} = 0.0582\,\mathrm{dex}$, median relative error $2.61\%$ above $0.01\times$ median signal, $x_{\mathrm{HI}}$ RMSE $= 0.00352$), outperforming $\lambda=0.3$ ($0.0584\,\mathrm{dex}$, $0.00478$), $\lambda=2.0$ ($0.0603\,\mathrm{dex}$, $0.00360$), $\lambda=4.0$ ($0.0569\,\mathrm{dex}$ with degraded median $2.82\%$), and dynamic auto-weighting ($0.0572\,\mathrm{dex}$, median $2.91\%$), whereas detaching gradients (\texttt{stop\_grad}) caused catastrophic $x_{\mathrm{HI}}$ failure ($\mathrm{RMSE}=0.4717$). Expanding the $\lambda=1.0$ joint architecture from three to eight random seeds yields the production release \textsc{PSX-emu v1.0} (\texttt{psx\_emu\_v1\_8seed}).

Before deploying \textsc{PSX-emu v1.0}, a provenance audit of the legacy 64-box development split (\texttt{val\_64}) uncovered exact parameter duplication: $54/64$ points matched rows inside the 100,000-point training table at zero standardized nearest-neighbor distance ($d_{\mathrm{NN}} = 0.0$), leaving only 10 genuinely disjoint points ($d_{\mathrm{NN}} \in [0.770, 1.480]$). To establish strict out-of-sample generalization and measure proximity to the intrinsic two-IC cosmic-variance floor, we generated an independent 300-point two-IC test suite stratified into three groups of 100 families: high-probability UV LF posterior draws (\texttt{LF\_draws}), active reionization draws (\texttt{selected\_active}), and 8D prior-edge points (\texttt{prior\_edges\_8D}). Nearest-neighbor verification against all 100,000 training inputs confirms zero overlaps ($0/300$ at $d_{\mathrm{NN}} < 10^{-5}$; overall $d_{\mathrm{NN}} \in [0.561, 1.400]$, median $1.034$).

\figwide{figures/fig6_psx_independent_audit.pdf}{Independent 300-point zero-overlap evaluation of the joint \textsc{PSX-emu v1.0} surrogate. (a) Power-spectrum $\log_{10}(\Delta^2_{21}+0.01)$ RMSE (dex) contrasting the contaminated 64-point split ($54$ overlapping vs.\ $10$ disjoint points, shaded top) against 300 strictly disjoint independent points ($d_{\mathrm{NN}} \ge 0.561$) across \texttt{emu\_joint\_l10} (3-seed, amber squares), \textsc{PSX-emu v1.0} (8-seed, navy diamonds), and the intrinsic two-IC half-difference noise floor (open circles). (b) Normalized multi-metric accuracy across all 300 independent points in percent for $\Delta^2_{21}>1\,\mathrm{mK}^2$, $\Delta^2_{21}>10\,\mathrm{mK}^2$, and Planck-normalized optical depth $\tau_e/\sigma_{\rm Planck}$ ($\sigma_{\rm Planck}=0.007$).}{fig:psx_audit}

Figure~\ref{fig:psx_audit} and Table~\ref{tab:psx_audit} summarize the 300-point zero-overlap audit. Expanding from 3 to 8 seeds reduces overall independent $\Delta^2_{21}$ RMSE from $0.0505$ to $0.0485\,\mathrm{dex}$ (less than twice the $0.0250\,\mathrm{dex}$ two-IC half-difference noise floor), reaching $0.0334\,\mathrm{dex}$ on \texttt{LF\_draws} (vs.\ $0.0276\,\mathrm{dex}$ floor) and $0.0382\,\mathrm{dex}$ on \texttt{selected\_active} (vs.\ $0.0238\,\mathrm{dex}$). At detectable amplitudes $\Delta^2_{21}>10\,\mathrm{mK}^2$, \textsc{PSX-emu v1.0} achieves $2.47\%$ median relative error ($7.52\%$ at $p_{90}$, $10.09\%$ at $p_{95}$ vs.\ $1.30\% / 4.53\% / 6.13\%$ intrinsic two-IC scatter), alongside raw (clipped) $x_{\mathrm{HI}}$ RMSE $= 0.00257$ ($0.00254$) and $\tau_e$ RMSE $= 1.72\times 10^{-4}$ ($2.45\%$ of $\sigma_{\rm Planck}=0.007$, $p_{90} = 3.06\%$ of $\sigma_{\rm Planck}$). Four-way power-law convergence diagnostics run for $20{,}000$ production steps ($\max \widehat{R} = 1.002577$) combining this verified eight-seed history head with original frozen PS models under Park \textit{et al.} noise plus a $20\%$ PS theory floor yield stable coeval/self-mock $t_\star$ medians ($0.504$, $0.535$, $0.600$, and $0.583$ across \texttt{LF}, \texttt{LF+}$\tau$\texttt{+}$x_{\mathrm{HI}}$, \texttt{PS}, and \texttt{PS+LF}), establishing sampler mixing prior to native lightcone truth-alignment tests.

\begin{table*}[!tb]
\caption{Zero-overlap independent 300-point audit of \texttt{emu\_joint\_l10} (3-seed) and \textsc{PSX-emu v1.0} (8-seed) vs.\ the two-IC half-difference noise floor ($N=300$ disjoint points; top two rows show the legacy 64-point split overlap diagnostic). Optical-depth ratios use $\sigma_{\rm Planck}=0.007$.}
\label{tab:psx_audit}
\scriptsize
\setlength{\tabcolsep}{2.2pt}
\begin{ruledtabular}
\begin{tabular}{lcccccc}
Subset / Model & PS RMSE ($\mathrm{dex}$) & $>1\,\mathrm{mK}^2$ med/$p_{90}$ & $>10\,\mathrm{mK}^2$ med/$p_{90}$ & $x_{\mathrm{HI}}$ RMSE & $\tau_e$ RMSE ($10^{-4}$) & $\tau_e$ $p_{90}/\sigma_{\rm P}$ \\
\hline
\texttt{val\_64} overlap 54 (3s / 8s) & 0.0581 / 0.0470 & 3.56\% / 2.91\% & 3.23\% / 2.62\% & 0.00292 / 0.00256 & 2.39 / 1.96 & 5.15\% / 4.80\% \\
\texttt{val\_64} disjoint 10 (3s / 8s) & 0.0586 / 0.0615 & 5.07\% / 3.62\% & 4.53\% / 3.21\% & 0.00454 / 0.00399 & 5.82 / 5.35 & 17.31\% / 9.46\% \\
\hline
Indep.\ 300, 3-seed (\texttt{l10}) & 0.0505 & 3.59\% / 10.78\% & 3.12\% / 8.98\% & 0.00305 & 1.96 & 3.22\% \\
Indep.\ 300, 8-seed (\textsc{PSX-emu}) & \textbf{0.0485} & \textbf{2.96\% / 9.40\%} & \textbf{2.47\% / 7.52\%} & \textbf{0.00254} & \textbf{1.72} & \textbf{3.06\%} \\
Indep.\ 300, 2-IC half-diff floor & 0.0250 & 1.53\% / 5.66\% & 1.30\% / 4.53\% & 0.00115 & 0.65 & 1.62\% \\
\hline
\texttt{LF\_draws} (100, 8s / 2-IC) & 0.0334 / 0.0276 & 2.75\% / 8.48\% & 2.35\% / 6.95\% & 0.00178 / 0.00104 & 1.21 / 0.68 & 3.08\% / 1.74\% \\
\texttt{selected\_active} (100, 8s / 2-IC) & 0.0382 / 0.0238 & 2.90\% / 9.29\% & 2.34\% / 6.94\% & 0.00284 / 0.00138 & 1.17 / 0.69 & 2.81\% / 1.57\% \\
\texttt{prior\_edges\_8D} (100, 8s / 2-IC) & 0.0669 / 0.0234 & 3.27\% / 10.50\% & 2.79\% / 8.77\% & 0.00286 / 0.00098 & 2.45 / 0.59 & 3.39\% / 1.54\% \\
\end{tabular}
\end{ruledtabular}
\end{table*}

\subsection{Redshift-Cadence Diagnosis and Localized RBF Residual Patching}\label{sec:plan_a}

When testing end-to-end PS+LF posterior inference against native \textsc{21cmFAST} mock spectra at benchmark truth ($t_\star = 0.500$, all other parameters at Park \textit{et al.} fiducial values), we diagnosed a discrete redshift-stepping mismatch between the coarse training grid ($\texttt{ZPRIME\_STEP\_FACTOR} = 1.02$) and dense native lightcone steps ($\texttt{1.01}$ and $\texttt{1.005}$). Evaluated at $t_\star=0.805$ versus $0.500$ across twelve cadence checks, $\chi^2$ remains $16.05\text{--}19.50$ on native dense grids ($\texttt{step}=1.02, 1.01, 1.005$) but collapses to $0.48$ once coarsened back to the emulator cadence. Directly sampling the frozen 100k power-spectrum base model (\texttt{emu\_conv100k.pt}) against native $\texttt{step}=1.005$ PS+LF mock data shifts the recovered $t_\star$ median upward to $0.6326$ ($0.524\sigma$ bias), with Fisher-propagated parameter bias $L_3(t_\star) = 0.971$ ($\max L_3 = 2.163$) and normalized spectrum error $L_2\ p_{90} = 0.7796$ at truth.

Rather than regenerating all 100,000 global simulations on a finer redshift grid, we construct a localized radial-basis-function (RBF) residual patch added purely on top of the untouched frozen 100k base predictions:
\begin{equation}
\widehat{y}_{\rm patched}(\theta) = \widehat{y}_{\rm base}(\theta) + \Delta\widehat{y}_{\rm RBF}(\theta),
\end{equation}
fitted via leave-one-out cross-validation strictly on training families drawn from the target posterior region. After two 24-training-family pilots ($\texttt{step}=1.01$, $8$ held-out families; Stage 1 PS-only and Stage 2 PS+LF in Table~\ref{tab:plan_a}), our scaled \textbf{Plan A} draws 80 disjoint families ($64$ training / $16$ held-out stratified across $16$ uniform $t_\star \in [0.05, 1.0]$ bins from the frozen joint posterior) $\times 2$ seeds $= 160$ full-resolution $512/128/250\,\mathrm{Mpc}$ simulations at $\texttt{ZPRIME\_STEP\_FACTOR}=1.005$ using GPU-accelerated \textsc{21cmFAST}. Cross-backend parity on $\texttt{task\_000}$ matches CHPC CPU execution to $1.94\times 10^{-16}$ median relative difference in coeval PS (all $361$ redshift nodes identical) and $2.22\times 10^{-11}$ in lightcone PS, while training-only LOO selects RBF hyperparameters $(\text{length}=2.0, \text{nugget}=0.001)$ without held-out tuning.

\figwide{figures/fig7_tstar_local_correction.pdf}{Localized RBF residual patching for redshift-cadence mismatch and $t_\star$ posterior recovery. (a) Reduction in $L_2$ $p_{90}$ spectrum residuals and $t_\star$ Fisher $L_3$ bias between frozen 100k base predictions (open amber) and RBF-patched predictions (filled navy) across Stages 1--3 relative to the $0.10$ accuracy floor (dotted). (b) Marginal $t_\star$ posterior medians and 68\% credible intervals relative to truth $t_\star=0.500$ (dashed), comparing four-way converged PL diagnostics (top four gray rows, $\max \widehat{R}=1.0026$) against Stage 2 (24-point) and Stage 3 Plan A (64-point, $\texttt{step}=1.005$) GPU MCMC recovery, where the $0.6326$ ($0.524\sigma$) base bias narrows to $0.4942$ ($0.017\sigma$).}{fig:plan_a}

Across all diagnostics (Fig.~\ref{fig:plan_a} and Table~\ref{tab:plan_a}), the 64-point Plan A patch reduces $L_2\ p_{90}$ by more than an order of magnitude at both truth ($0.7796 \to 0.0622$) and across 16 disjoint held-out families ($0.7924 \to 0.0671$), drives held-out median $t_\star$ Fisher $L_3$ from $0.973$ down to $0.052$ ($13/16 \le 0.10$; truth $t_\star$ $L_3 = 0.971 \to 0.106$), and restores full $64\times 3000$-step PS+LF GPU MCMC recovery of $t_\star$ from $0.6326$ ($0.524\sigma$) to $0.4942$ ($0.017\sigma$; 68\% interval $[0.1772, 0.8393]$). We also transparently document remaining worst-case multi-parameter bounds: holding all 8 coordinates strictly below $L_3 \le 0.10$ is not yet uniformly met ($\max L_3 = 0.440$ at truth from $\alpha_{\rm esc}$, vs.\ $2.163$ pre-patch; held-out all-parameter $\max L_3$ median $0.303$, worst-case marginal family $1.780$ vs.\ $2.973$ pre-patch), defining a clear target for covariance-aware anisotropic Gaussian-process patching in future work.

\begin{table*}[!tb]
\caption{Progressive localized RBF residual patch verification across Stage 1 (24pt PS-only, $\texttt{step}=1.01$), Stage 2 (24pt PS+LF, $\texttt{step}=1.01$), and Stage 3 Plan A (64pt PS+LF, $\texttt{step}=1.005$, 160 GPU runs). Parenthetical values give pre-patch frozen base metrics.}
\label{tab:plan_a}
\scriptsize
\setlength{\tabcolsep}{2.2pt}
\begin{ruledtabular}
\begin{tabular}{lccccc}
Stage / Configuration & Truth $L_2\ p_{90}$ & Held $L_2\ p_{90}$ & Truth $t_\star$ (/ max) $L_3$ & Held med $t_\star\ L_3$ & MCMC $t_\star$ (bias) \\
\hline
Stage 1: 24pt PS-only (8 held) & 0.0406 (0.7608) & 0.1126 (0.7829) & 0.161 / 0.227 (0.763 / 1.463) & \dots & 0.5356 / $0.123\sigma$ ($0.500\sigma$) \\
Stage 2: 24pt PS+LF (8 held) & 0.0406 (0.7608) & 0.1126 (0.7829) & 0.223 / 0.223 (0.920 / 2.029) & 0.225 (0.922) & 0.6038 / $0.363\sigma$ ($0.534\sigma$) \\
Stage 3: 64pt Plan A (16 held) & \textbf{0.0622} (0.7796) & \textbf{0.0671} (0.7924) & \textbf{0.106} / \textbf{0.440} (0.971 / 2.163) & \textbf{0.052} (0.973) & \textbf{0.4942 / $0.017\sigma$} ($0.524\sigma$) \\
\end{tabular}
\end{ruledtabular}
\end{table*}

\subsection{Anisotropic Mat\'ern-5/2 ARD-GP Lightcone Patches and Three-Way 8D Hybrid Stretch+GMM Chains}\label{sec:ard_gmm}

To generalize localized residual patching beyond isotropic coeval RBF interpolation to full 21-cm lightcone spectra and broken-power-law backgrounds without retraining \textsc{PSX-emu v1.0}, we fit four anisotropic Mat\'ern-5/2 automatic-relevance-determination Gaussian-process (ARD-GP) patches across independent 64-run native two-IC training sets ($\texttt{step}=1.005$, $512/128/250\,\mathrm{Mpc}$) and evaluate each on 16 disjoint held-out families plus dedicated native truth runs (Fig.~\ref{fig:ard_gmm}a and Table~\ref{tab:ard_gmm}). Across all four targets, ARD-GP patching compresses held-out normalized $L_2\ p_{90}$ residuals by more than an order of magnitude: on $\mathrm{PL}\,(10, 0.968)$ coeval (\texttt{pl\_kp10\_ms0p968\_coeval\_ard\_global.npz}) from $0.7924$ to $0.0608$ (nugget $1.51\times 10^{-4}$; \texttt{step\_1005} $t_\star\ L_3 = 0.100$); on $\mathrm{PL}\,(10, 0.968)$ lightcone (\texttt{pl\_kp10\_ms0p968\_lc\_ard\_global.npz}) from $3.9249$ to $0.0656$ (LOO MSE $0.17043$, nugget $3.53\times 10^{-5}$); on $\mathrm{BT}\,(10, 2.5)$ lightcone (\texttt{bt\_kp10\_ms2p5\_lc\_ard\_global.npz}) from $1.2988$ to $0.0278$ ($1.2005 \to 0.0124$ at truth \texttt{task\_160/161}, nugget $3.29\times 10^{-4}$); and on $\mathrm{BT}\,(3, 2.0)$ lightcone (\texttt{bt\_kp3\_ms2p0\_lc\_ard\_global.npz}) from $0.6699$ to $0.0552$ ($0.8043 \to 0.0921$ at truth, LOO MSE $0.42756$, nugget $8.48\times 10^{-4}$).

\figwide{figures/fig8_ard_gp_and_threeway_gmm.pdf}{Anisotropic Mat\'ern-5/2 ARD-GP local lightcone/coeval patching and three-way 8D hybrid Stretch+GMM posterior accuracy. (a) Pre-patch frozen base (open amber) vs.\ post-patch ARD-GP (filled navy) normalized residual $L_2\ p_{90}$ relative to the $0.10$ floor (dotted) across all four local patches (16 held-out families and native truth boxes). (b) Standardized parameter recovery $(q - \theta_{\mathrm{truth}})/\sigma_{68}$ (median and 68\% credible intervals vs.\ $\pm 1\sigma$ shaded band) from $230{,}400$ post-burn samples per model for $\mathrm{PL}\,(10,0.968)$ (navy diamonds), $\mathrm{BT}\,(10,2.5)$ (amber circles), and $\mathrm{BT}\,(3,2.0)$ (teal squares).}{fig:ard_gmm}

Pairwise 8D posterior sampling combines $50\%$ affine-invariant stretch moves with $50\%$ four-component Gaussian-mixture-model (4-GMM) independence proposals across two independent ensembles of $64$ walkers run for $20{,}000$ steps ($2{,}000$ burn-in, thin $10$, yielding $230{,}400$ post-burn samples per benchmark; Table~\ref{tab:ard_gmm} and Fig.~\ref{fig:ard_gmm}b). The baseline $\mathrm{PL}\,(10, 0.968)$ chain passes all strict multi-ensemble convergence criteria (\texttt{convergence\_established=True}, $\max\widehat{R}=1.0071$, $\min[N/\tau_{\mathrm{ens}}]=141.0$, $\min[N/\tau_{\mathrm{walk}}]=181.7$, ensemble shift $0.0201\sigma$) and recovers all eight parameters within $|\mathrm{med\_z}|\le 0.383$, including $t_\star = 0.4850^{+0.3493}_{-0.2882}$ ($0.047\sigma$), $L_X = 40.5098$, and $E_0 = 500.1\,\mathrm{eV}$. Under stronger spectral enhancement at $\mathrm{BT}\,(3, 2.0)$ around an LF-matched broken truth ($f_{\star,10}=-1.7491$, $\alpha_\star=0.9077$, $M_{\rm turn}=8.0100$, $t_\star=0.5045$), the chain also passes strict convergence ($\max\widehat{R}=1.0049$, $\min[N/\tau_{\mathrm{ens}}]=141.1$, $\min[N/\tau_{\mathrm{walk}}]=275.7$), closely recovering star-formation slope and efficiency ($f_{\star,10}=-1.6878$, $+0.275\sigma$; $\alpha_\star=0.8886$, $-0.299\sigma$) while $M_{\rm turn}=8.3996$ ($+1.351\sigma$) and $L_X=41.2638$ ($+1.489\sigma$) exhibit mock-noise/prior-volume pull. At $\mathrm{BT}\,(10, 2.5)$ ($f_{\star,10}=-1.2545$, $\alpha_\star=0.5859$ at $+1.403\sigma$, $M_{\rm turn}=8.5091$), $\max\widehat{R}$ reaches $1.0105$ in $L_X$ ($0.0005$ above $1.01$, $\min[N/\tau_{\mathrm{ens}}]=105.3$, 4-GMM acceptance $5.8\text{--}6.1\%$), and we flag \texttt{convergence\_established=False} rather than asserting full convergence. Note that large shifts of $\mathrm{BT}\,(3, 2.0)$ relative to the $\mathrm{PL}$ chain ($+5.454\sigma_{\mathrm{PL}}$ in $\alpha_\star$ and $+6.762\sigma_{\mathrm{PL}}$ in $L_X$) reflect the altered UV-LF-matched broken truth coordinates ($f_{\star,10}, \alpha_\star, M_{\rm turn}$) required to preserve fixed galaxy luminosity functions under enhanced low-mass halo abundance.

\begin{table*}[!tb]
\caption{Summary of the four Mat\'ern-5/2 ARD-GP local patches (64 train / 16 held-out, $\texttt{step}=1.005$) and three-way 8D hybrid Stretch+GMM posterior chains ($2\times 64\times 20{,}000$, burn-in $2{,}000$, thin $10$, $230{,}400$ post-burn samples).}
\label{tab:ard_gmm}
\scriptsize
\setlength{\tabcolsep}{1.8pt}
\begin{ruledtabular}
\begin{tabular}{lccccc}
Model / Patch & Held $L_2\ p_{90}$ (Pre$\to$Post) & Truth $L_2\ p_{90}$ & $\max \widehat{R}$ & $\min(N/\tau_{\mathrm{ens}})$ / Shift & Key recovery / Flag \\
\hline
$\mathrm{PL}\,(10, 0.968)$ coeval & $0.7924 \to 0.0608$ & \dots & \dots & \dots & \texttt{step\_1005} $t_\star\ L_3=0.100$ \\
$\mathrm{PL}\,(10, 0.968)$ LC + 8D & $3.9249 \to 0.0656$ & \dots & 1.0071 & 141.0 / $0.020\sigma$ & \textbf{PASS}; $|\mathrm{med\_z}|\le 0.383$ \\
$\mathrm{BT}\,(10, 2.5)$ LC + 8D & $1.2988 \to 0.0278$ & $1.2005 \to 0.0124$ & 1.0105 & 105.3 / $0.031\sigma$ & \textbf{FLAG}: $L_X$ $\widehat{R}=1.0105$, $\alpha_\star +1.40\sigma$ \\
$\mathrm{BT}\,(3, 2.0)$ LC + 8D & $0.6699 \to 0.0552$ & $0.8043 \to 0.0921$ & 1.0049 & 141.1 / $0.014\sigma$ & \textbf{PASS}; $L_X +1.49\sigma$, $M_{\rm turn} +1.35\sigma$ \\
\end{tabular}
\end{ruledtabular}
\end{table*}

\section{Discussion and Outlook}\label{sec:discussion}

Having established and verified the 10D history emulator, \textsc{PSX-emu v1.0}, Plan A localized RBF patching, and four anisotropic Mat\'ern-5/2 ARD-GP lightcone/coeval patches across prior splits, 300 disjoint test points, and three-way 8D hybrid Stretch+GMM benchmarks, three methodological boundaries define the roadmap for downstream multi-observable joint inference:

\begin{enumerate}
\item \textbf{Domain hygiene, anisotropic residual GP patching, and IC variance}: Our pilots and \texttt{val\_64} overlap audit show that full-domain $m_s \in [0.5, 4.0]$ full-box spherical-FFT labels must always be paired with strict zero-overlap verification ($d_{\mathrm{NN}}>0$). Section~\ref{sec:ard_gmm} demonstrates that replacing isotropic RBF kernels with anisotropic Mat\'ern-5/2 ARD-GP patches suppresses held-out lightcone and coeval $L_2\ p_{90}$ residuals to $0.0278\text{--}0.0656$ without retraining base weights; propagating full GP predictive covariance into the PS+LF likelihood alongside generative multi-IC uncertainty models \cite{Breitman2026} offers a direct path to bound remaining mock-degeneracy pulls ($M_{\rm turn}$, $L_X$, $\alpha_\star$).
\item \textbf{Multi-ensemble sampler convergence and mock-only screening}: High-dimensional GPU-accelerated scans must enforce multi-ensemble split-$\widehat{R}$ diagnostics \cite{Vehtari2021} and parameter-truth alignment before claiming inference acceptance. As a concrete cautionary control, an exploratory 1,000,000-step 6D history-emulator GPU chain (\texttt{mock\_only}, single ensemble per $(k_{p,h}, m_s)$ cell, $12.40\%$ acceptance, $03\mathrm{h}17\mathrm{m}$ runtime) failed global convergence screening (\texttt{all\_local\_diagnostics\_pass=false}, missing multi-chain $\widehat{R}$, misaligned X-ray constants) and was rejected from accepted physical results.
\item \textbf{End-to-end runtime accounting}: While a five-seed history-emulator ensemble evaluates in $\sim 0.41\,\mathrm{ms}$ per step versus $416\,\mathrm{s}$ median wall time for a single 16-thread direct \textsc{21cmFAST} box---and GPU Wave-6 execution accelerates full $512/128/250\,\mathrm{Mpc}$ $\texttt{step}=1.005$ verification runs to $\sim 3.5\,\mathrm{min}$ per box---end-to-end inference chain timing continues to include exact UV LF calculation and optical-depth integration.
\end{enumerate}

\section{Conclusions and Data Availability}\label{sec:conclusions}

We have established the architecture, simulation contracts, random-IC likelihood semantics, full-box power-spectrum production, independent multi-task validation, and localized residual patching of the ten-dimensional reionization emulation framework for broken-power-law small-scale matter spectra in \textsc{21cmFAST}. Six principal methodological conclusions emerge:
\begin{enumerate}
\item \emph{Ensembling and coverage in 10D history emulation}: Five-seed physical averaging reduces median trajectory RMSE by $21\text{--}34\%$ over individual seeds, and expanding single-IC scrambled Sobol coverage from 65,536 to 100,000 runs improves $x_{\mathrm{HI}}(5.9)$ RMSE by $35.6\%$, $\tau$ RMSE by $23.2\%$, and prior-box likelihood tails by $\sim 26\%$ across paired development splits.
\item \emph{Thread-pinned GSL reproducibility}: Because native \textsc{21cmFAST} partitions GSL Fourier-mode random streams by OpenMP thread count, exact multi-IC pairing and backend verification strictly require pinning \texttt{N\_THREADS=16} alongside integer random seeds and binary SHA-256 digests.
\item \emph{Reconciliation of random-IC likelihood semantics}: Although nonlinear threshold amplification causes large IC scatter in prior-box extremes ($q_{90}[s_{\mathrm{IC}}(\ell_{\mathrm{hist}})] = 0.948$), direct multi-IC and stratified out-of-sample simulations prove that within the observational posterior support the 100k history emulator achieves $q_{90}(|\Delta\ell_{\mathrm{reion}}|) \le 0.056$ and tracks four-IC marginal likelihoods within $\le 0.028$ ($\ell(\bar{X}_4)$ within $\le 0.0054$; $|\ell_{\mathrm{emu}}-\ell(\bar{X}_4)| \le 0.033$).
\item \emph{Full-box data engineering and 300-point \textsc{PSX-emu v1.0} validation}: Controlled pilots show single-IC coverage dominates multi-IC replication while 12-subbox averaging introduces up to $1.48\times$ parameter-dependent low-$k$ bias. On the completed 100,000-box \texttt{ps\_xhi\_fullbox\_v1} dataset, the eight-seed multi-task \textsc{PSX-emu v1.0} surrogate ($\lambda=1.0$) reaches $0.0485\,\mathrm{dex}$ overall PS RMSE ($0.0334\,\mathrm{dex}$ on LF draws; $2.47\%$ median relative error at $>10\,\mathrm{mK}^2$) and $2.45\%$ $\tau_e/\sigma_{\rm Planck}$ RMSE across 300 zero-overlap points.
\item \emph{Localized RBF patching for cadence mismatch}: A 64-family GPU-produced Plan A RBF patch ($160$ $\texttt{step}=1.005$ boxes) reduces held-out $L_2\ p_{90}$ residuals from $0.7924$ to $0.0671$ and held-out median $t_\star$ Fisher $L_3$ from $0.973$ to $0.052$, removing frozen-base $t_\star$ PS+LF MCMC bias ($0.6326 \to 0.4942$, $0.524\sigma \to 0.017\sigma$) without retraining base weights.
\item \emph{Anisotropic ARD-GP lightcone patches and three-way 8D Stretch+GMM convergence}: Across four 64-run Mat\'ern-5/2 ARD-GP benchmarks, held-out $L_2\ p_{90}$ residuals drop to $0.0278\text{--}0.0656$. Three-way $230{,}400$-sample hybrid Stretch+GMM chains satisfy strict multi-ensemble convergence at $\mathrm{PL}\,(10,0.968)$ ($\max\widehat{R}=1.0071$, $|\mathrm{med\_z}|\le 0.383$) and $\mathrm{BT}\,(3,2.0)$ ($\max\widehat{R}=1.0049$) while flagging marginal $\widehat{R}=1.0105$ in $L_X$ at $\mathrm{BT}\,(10,2.5)$.
\end{enumerate}

Contracts, training configurations, validation JSON/NPZ receipts, and figure scripts for the history and joint power-spectrum emulators are maintained at \url{https://github.com/HainaHuang729/21cm-reionization-emulator}. Companion physical parameter and small-scale matter-power inference results obtained with the 100k history emulator and exact UV LFs are hosted at \url{https://github.com/HainaHuang729/early-galaxies-reionization-bt}, and interactive 3D/2D/1D visualization views of the underlying 21-cm simulations are provided at \url{https://github.com/HainaHuang729/21cm-reionization-3d}.

\begin{acknowledgments}
Computations reported in this work were carried out on the Central High-Performance Computing (CHPC) Cluster and Department of Physics GPU computing facilities at The Chinese University of Hong Kong (CUHK). We thank the CUHK HPC support staff for assistance with job accounting and node scheduling across large-scale \textsc{21cmFAST} simulation sweeps.
\end{acknowledgments}

\appendix

\section{Contract Digests and Grid Specifications}\label{app:provenance}

Table~\ref{tab:provenance} records the frozen SHA-256 checksums, coordinate grids, and numerical settings linking the 10D history emulator, the \texttt{T7}/\texttt{A3} IC audits, and the \texttt{ps\_xhi\_fullbox\_v1} dataset for deterministic reproduction.

\begin{table}[!tbp]
\caption{Frozen numerical contracts and coordinate specifications.}
\label{tab:provenance}
\footnotesize
\setlength{\tabcolsep}{3.5pt}
\begin{ruledtabular}
\begin{tabular}{p{3.05cm}p{4.83cm}}
Contract / grid item & Exact specification \\
\hline
Native binary SHA-256 & \texttt{3211a6629109da73...94e5702b2} \\
Simulation box & $\texttt{BOX\_LEN}=250\,\mathrm{Mpc}$, $\texttt{DIM}=512$, $\texttt{HII\_DIM}=128$ \\
Threading & \texttt{N\_THREADS=16} (pinned GSL RNG stream) \\
History redshifts (32) & $\{5.0, 5.9, 6, 7, \dots, 35\}$ \\
$\tau$ adapter grid (31) & Uniform $z \in [5, 35]$, 31 interpolation nodes \\
Native PS shells (60) & Uniform $|k| \in [0.0232, 2.7628]\,\mathrm{Mpc}^{-1}$ \\
Target PS centers (32) & Log-spaced $k \in [0.1000, 1.0000]\,\mathrm{Mpc}^{-1}$ \\
PS loss / eval floor & $\log_{10}(\Delta^2_{21}/\mathrm{mK}^2 + 0.01)$ \\
T7 audit acceptance SHA & \texttt{090c4f6fdfb08cee...3c109738} \\
65k / 100k history-emulator SHA & \texttt{ac53130045034015} / \texttt{371f110162ed9ef3} \\
\texttt{emu\_joint\_l10.pt} SHA & \texttt{1ed7e13c036d7322...451881faf} \\
\texttt{psx\_emu\_v1\_8seed.pt} SHA & \texttt{c9623d018340b3ef...09dd23059} \\
Plan A RBF patch & 64 train / 16 held, $\texttt{step}=1.005$ ($l=2.0,\,\eta=10^{-3}$) \\
PL coeval / LC ARD-GP SHA & \texttt{694ea54d30388f61} / \texttt{17e05f6037c86e1a} \\
BT10 / BT3 LC ARD-GP SHA & \texttt{70756f67701e9df8} / \texttt{6db69c21f984c932} \\
PL / BT10 / BT3 GMM SHA & \texttt{6475e4ba} / \texttt{ac8446ca} / \texttt{59792a8a} \\
\end{tabular}
\end{ruledtabular}
\end{table}

\bibliography{emulator_methods}

\end{document}
